
\subsection{洛伦兹力}

\begin{tabularx}{\columnwidth}{XX}
  洛伦兹力&$F=qvB\sin\theta$\\
  向心力关系&$qvB=m\dfrac{v^2}{R}$\\
  半径&$R=\dfrac{mv}{qB}$\\
  周期&$T=\dfrac{2\pi m}{qB}$\\
  角速度&$\omega=\dfrac{qB}{m}$\\
  运动时间&$t=\dfrac{\theta m}{qB}=\dfrac{\theta R}{v}$\\
\end{tabularx}

\begin{formulanotebox}
洛伦兹力始终垂直于速度方向，因此单独作用时不做功，只改变速度方向，不改变速率。

公式中的 $q$ 常用电荷量大小；判断偏转方向时必须结合电荷正负。负电荷受力方向与正电荷相反。
\end{formulanotebox}

\begin{keypointbox}[title={\strut\textcolor{white}{\bfseries 重点知识点：圆周偏转模型}}]
带电粒子垂直进入匀强磁场后做匀速圆周运动。解题主线是：
\[
  \text{定圆心} \rightarrow \text{找半径} \rightarrow \text{求圆心角} \rightarrow \text{算时间}.
\]
轨迹圆心一定在入射点速度方向的垂线上，速度偏转角等于轨迹圆对应的圆心角。
\end{keypointbox}

\begin{casetypebox}[title={\strut\textcolor{white}{\bfseries 常见题型：边界磁场}}]
\topicimage[0.58\columnwidth]{assets/直线边界，出发点在边界上.png}
\topicimage[0.58\columnwidth]{assets/直线边界，出发点在磁场中.png}
\begin{itemize}[leftmargin=1.5em]
  \item 直线边界中，粒子从边界进入又从同一边界离开时，进入夹角等于离开夹角。
  \item 两平行边界间距为 $d$ 时，可用几何关系 $d=R(1-\cos\theta)$ 或弦长关系求出射位置。
  \item 圆形磁场中若粒子沿半径进入，常沿另一条半径方向离开。
  \item 速度可变时，把速度变化转化为半径变化，临界轨迹通常与边界或顶点相切。
\end{itemize}
\end{casetypebox}

\begin{casetypebox}[title={\strut\textcolor{white}{\bfseries 常见题型：磁聚焦}}]
同速同荷质比粒子在同一匀强磁场中半径相同。磁聚焦问题不要逐条轨迹硬算，优先寻找入射方向、出射方向和轨迹圆心的对称关系。

\begin{center}
  \begin{tikzpicture}[scale=0.78, >=Stealth, line cap=round, line join=round]
  \begin{scope}
    \draw[thick, fill=blue!5] (0,0) circle (1.8);
    \node at (0,-2.25) {\small 圆形磁场聚焦};
    \node[blue!65!black] at (0,1.45) {\small $B$};

    \coordinate (A) at (-1.8,0);
    \coordinate (F) at (1.8,0);
    \coordinate (C1) at (0,1.8);
    \coordinate (C2) at (0,0.95);
    \coordinate (C3) at (0,0.25);

    \draw[->, thick] (-2.55,1.15) -- (-1.82,1.15);
    \draw[->, thick] (-2.55,0.00) -- (-1.82,0.00);
    \draw[->, thick] (-2.55,-1.15) -- (-1.82,-1.15);

    \draw[orange!85!black, very thick, ->] (-1.38,1.15)
      .. controls (-0.45,1.65) and (0.95,1.15) .. (1.45,0.15);
    \draw[orange!85!black, very thick, ->] (A)
      .. controls (-0.55,1.25) and (0.65,1.25) .. (F);
    \draw[orange!85!black, very thick, ->] (-1.38,-1.15)
      .. controls (-0.15,-0.15) and (0.9,0.05) .. (1.45,-0.15);

    \fill (F) circle (1.6pt);
    \node[right] at (F) {\small 焦点};
  \end{scope}

  \begin{scope}[xshift=6.1cm]
    \draw[thick, fill=blue!5] (-1.2,0) arc (180:90:1.2) -- (1.2,1.2)
      arc (90:0:1.2) -- cycle;
    \node at (0,-2.25) {\small 叶形磁场聚焦};
    \node[blue!65!black] at (0,1.35) {\small $B$};

    \coordinate (P) at (-1.2,0);
    \coordinate (M) at (0,1.2);
    \coordinate (Q) at (1.2,0);

    \draw[->, thick] (-2.2,0) -- (P);
    \draw[orange!85!black, very thick, ->] (P) arc (180:90:1.2);
    \draw[orange!85!black, very thick, ->] (M) arc (90:0:1.2);
    \draw[->, thick] (Q) -- (2.2,0);

    \fill (P) circle (1.5pt);
    \fill (M) circle (1.5pt);
    \fill (Q) circle (1.5pt);
    \node[below] at (P) {\small 入射};
    \node[above] at (M) {\small 转接};
    \node[below] at (Q) {\small 出射};
    \node at (-0.62,0.58) {\small $\frac14 T$};
    \node at (0.62,0.58) {\small $\frac14 T$};
  \end{scope}
\end{tikzpicture}

\end{center}
\end{casetypebox}
